% Original: PDF strana 11, gornji slajd; prvi deo.
\slajdid{03-s021}
\begin{frame}[t,label=03-s021]{Lanac kodera prioriteta sa 16 ulaza}
\setlength{\parskip}{0pt}
% element: 03-s021:e02
Sa ovako realizovanim koderom prioriteta mreža sa \textbf{16 ulaza} postaje:
\par\vspace{1mm}
\begin{columns}[c,onlytextwidth]
\begin{column}{.47\textwidth}\centering
% element: 03-s021:e01
\begin{tikzpicture}[circuit logic US,x=.82mm,y=.67mm,font=\fontsize{7.8}{9}\selectfont,
 zbir/.style={or gate US,draw,line width=.8pt,logic gate inputs=nnnn,minimum width=4.8mm,minimum height=5mm,scale=.52}]
\foreach \g/\yy in {3/54,2/36,1/18,0/0}{
 \begin{scope}[shift={(0,\yy)}]
 \draw[DEvod] (0,0) rectangle (16,16);
 \node at (8,9) {KP};
 \foreach \n/\yin in {3/12.5,2/9.4,1/6.3,0/3.2}{
  \pgfmathtruncatemacro{\ii}{4*\g+\n}
  \node[anchor=west,inner sep=.5pt] at (0,\yin) {I\n};
  \draw[DEvod] (-2,\yin)--(0,\yin);
  \node[anchor=east,inner sep=.5pt] at (-2,\yin) {I\ii};
 }
 \foreach \port/\yout/\lab in {gs/10.8/GS,a1/6.5/A1,a0/3.2/A0}{
  \node[anchor=east,inner sep=.5pt] at (16,\yout) {\lab};
  \draw[DEvod] (16,\yout)--(18,\yout);
  \coordinate (p\g-\port) at (18,\yout);
 }
 \node[anchor=north,inner sep=.5pt] at (8,16) {EI};
 \draw[DEvod] (8,16)--(8,17);
 \coordinate (p\g-ei) at (8,17);
 \node[anchor=south,inner sep=.5pt] at (8,0) {EO};
 \draw[DEvod] (8,0)--(8,-1);
 \coordinate (p\g-eo) at (8,-1);
 \end{scope}
}
\draw[DEvod,DEisticanje] (p3-ei)--++(-3,0) node[left,inner sep=.5pt] {1};
\foreach \g/\next in {3/2,2/1,1/0}
 \draw[DEvod,DEisticanje] (p\g-eo)--(p\next-ei);
\draw[DEvod] (47,27) rectangle (63,43);
\node at (55,36) {KP};
\foreach \n/\yin in {3/39.5,2/36.4,1/33.3,0/30.2}
 \node[anchor=west,inner sep=.5pt] at (47,\yin) {I\n};
\foreach \g/\trunk/\yin in {3/43/39.5,2/41/36.4,1/37/33.3,0/39/30.2}
 \draw[DEvod] (p\g-gs)--(\trunk,0 |- p\g-gs)|-(47,\yin);
\node[anchor=north,inner sep=.5pt] at (55,43) {EI};
\draw[DEvod,DEisticanje] (55,43)--(55,45) node[above,inner sep=.5pt] {1};
\foreach \port/\yy/\inner/\outer in {gs/37.8/GS/GS,a1/33.5/A1/A3,a0/30.2/A0/A2}{
 \node[anchor=east,inner sep=.5pt] at (63,\yy) {\inner};
 \draw[DEvod] (63,\yy)--(65,\yy) node[right,inner sep=.5pt] {\outer};
}
\node[anchor=south,inner sep=.5pt] at (55,27) {EO};
\draw[DEvod] (55,27)--(55,25)--(65,25) node[right,inner sep=.5pt] {EO};
\node[zbir] (a1) at (40,-1) {};
\node[zbir] (a0) at (40,-10.5) {};
\foreach \g/\pin/\xx in {3/1/31,2/2/27,1/3/23,0/4/19}
 \draw[DEvod] (p\g-a1)--(\xx,0 |- p\g-a1)|-(a1.input \pin);
\foreach \g/\pin/\xx in {3/1/33,2/2/29,1/3/25,0/4/21}
 \draw[DEvod] (p\g-a0)--(\xx,0 |- p\g-a0)|-(a0.input \pin);
\foreach \name/\lab in {a1/A1,a0/A0}
 \draw[DEvod] (\name.output)--++(2,0) node[right,inner sep=.5pt] {\lab};
\end{tikzpicture}

\end{column}
\begin{column}{.50\textwidth}
% element: 03-s021:e03
Mrežu možemo dodatno pojednostaviti: izlaze \(A_1,A_0\) neaktivnog kodera
postavimo u \textcolor{DEisticanje}{\textbf{stanje visoke impedanse}}.
\par\vspace{2mm}
Tada odgovarajuće izlaze možemo \textbf{direktno spojiti},
pa nam \textcolor{DEisticanje}{\textbf{četvoroulazni ILI gejtovi nisu potrebni}}.
\end{column}
\end{columns}
\end{frame}
